\chapter{绪论}
\label{chap:intro}

\section{研究背景与意义}

近年来，数字经济在国民经济中的占比持续提升，正深刻改变要素配置方式与产业组织形态。然而，由于资源禀赋与发展阶段不同，各地区数字经济发展呈现明显的不平衡特征\citep{xugc2023digital}。如何科学测度数字经济发展水平、客观刻画其区域差异，已成为经济统计领域的重要议题\citep{liwei2022region}。

本文的研究意义主要体现在两个方面。其一，理论上，构建一套多维度、可比较的数字经济发展水平测度框架，有助于丰富数字经济量化研究\footnote{本文为模板演示文档，全部数据与结论均为虚构，仅用于展示排版效果。}。其二，实践上，区域差异的来源分解可为差异化政策制定提供参考。

\section{研究内容与研究方法}

围绕“测度—比较—解释”这一主线，本文拟回答以下问题：

\begin{enumerate}[label=（\arabic*）]
    \item 如何从多个维度构建数字经济发展水平的综合指标？
    \item 样本期内数字经济发展水平的总体趋势如何？
    \item 区域差异主要来源于组内差距还是组间差距？
\end{enumerate}

为此，本文综合运用以下方法：采用熵值法客观确定指标权重，避免主观赋权的随意性；采用泰尔指数对区域差异进行可加分解；并辅以聚类分析对地区进行分组比较\citep{chen2019cluster}。

\section{研究框架}

本文共分五章。第~\ref{chap:method}~章给出指标体系与测度方法，总体技术路线见图~\ref{fig:roadmap}；第~\ref{chap:empirical}~章在虚构数据上展开实证分析；末章总结结论并讨论展望。
